%%%PAGE 232 rot=0 legibility=good summary=九条简短要点/易错提示：合力方向动能定理与动滑轮、波的叠加多解、干涉减弱点、克服功正负、变压器电路、空气中光传播时间、列函数对照图象、旧欧姆表、多解情形（页面下部为透印的淡字，已忽略）
%%%BLOCK id=232-01 ch=06 sec=动能定理·合力方向与动滑轮 kind=note alt=09,03 cont=none y=0.00-0.10
\textbf{1、}
%FIG p232_a 0.07 0.008 0.175 0.058
\notefig[0.2]{figures/p232_a.png}
$F_{\text{合}}$在\unclear{$F_{\text{合}}$}方向上用动能定理，动滑轮2倍力
%%%END
%%%BLOCK id=232-02 ch=08 sec=波的叠加·多解与最小公倍数 kind=problem alt=none cont=none y=0.12-0.235
\begin{problem}[2]
\begin{solution}
$x_{\text{甲}}=50+50n_1$ \qquad $V_{\text{甲}}=V_{\text{乙}}=25\unit{cm/s}$\par
$x_{\text{乙}}=50+60n_2$ $\;\to\;$ $x=50+300n$ \quad 找最小公倍数\par
波程差 $\Delta x=160n_2-50n_{\text{甲}}+5\ge 5$ % CHECK: 160 的首位"1"疑为多余笔画（按 x 乙 的式子似应为 60n_2）
\qquad $t_{\min}=\dfrac{\Delta x}{\text{\unclear{$\blacksquare$}}V_{\text{相}}}=0.15$
\end{solution}
\end{problem}
%%%END
%%%BLOCK id=232-03 ch=08 sec=波的干涉·减弱点 kind=problem alt=none cont=none y=0.25-0.33
\begin{problem}[3]
%FIG p232_b 0.07 0.262 0.20 0.30
\notefig[0.2]{figures/p232_b.png}
\begin{solution}
设$x$为\unclear{减弱}点
\[ x-(2-x)=(2n+1)\frac{\lambda}{2} \qquad x=\frac{1}{2}n+\frac{5}{4} \qquad x\in[0,2] \]
故 $n=-2,-1,0,1$
\end{solution}
\end{problem}
%%%END
%%%BLOCK id=232-04 ch=06 sec=功·克服功的正负 kind=note alt=none cont=none y=0.335-0.375
\textbf{4、}克服功为正、功可正负
%%%END
%%%BLOCK id=232-05 ch=13 sec=变压器·含变压器电路 kind=note alt=10 cont=none y=0.38-0.44
\textbf{5、}
%FIG p232_c 0.075 0.372 0.215 0.442
\notefig[0.2]{figures/p232_c.png}
$U_1=U_2+U_{\text{变}}$ （变压器电压） \qquad 电容器上$4$\unclear{V}为额定电压 % CHECK: 4 后的字母手写像 N，应为 V？
%%%END
%%%BLOCK id=232-06 ch=15 sec=光传播时间·空气另算 kind=note alt=11 cont=none y=0.445-0.525
\textbf{6、}
%FIG p232_d 0.095 0.438 0.164 0.52
\notefig[0.2]{figures/p232_d.png}
\begin{keybox}
空气时间另算 // 无\unclear{场}区匀速阶段
\end{keybox}
空气透镜 \quad $t=\text{\unclear{$\bullet$}}\dfrac{s}{c}$ \quad $n=1$ % CHECK: 分式前有一个辨认不清的小符号（像带竖线的圆点）
%%%END
%%%BLOCK id=232-07 ch=03 sec=动力学·列函数与图象对照 kind=method alt=90 cont=none y=0.53-0.62
\textbf{7、}列方程成函数与图象对照得数值
%FIG p232_e 0.465 0.535 0.595 0.62
\notefig[0.2]{figures/p232_e.png}
%%%END
%%%BLOCK id=232-08 ch=10 sec=多用电表·欧姆表(旧电池) kind=note exp=1 alt=none cont=none y=0.60-0.66
\textbf{8、}\unclear{旧}欧姆表用公式算 $R_{\text{测}}=R_{\text{真}}$\par
（直接测 $R_{\text{测}}$偏大）
%%%END
%%%BLOCK id=232-09 ch=07 sec=多解问题·易错点 kind=note alt=01,03 cont=none y=0.66-0.70
\textbf{9、}上抛、碰撞、$f$的两种方向可能性
%%%END
%%%PAGEEND 232
